\begin{figure*}[ht]
\centering
\includegraphics[width=0.95\linewidth]{fig/dataset_teaser_0318.png}
\caption{\textbf{{Overview of OmniViTac dataset.}} \textbf{Left:} The Cross-Embodiment Data Collection Platform features UNFactory 7-DoF xArm and the TacUMI manipulation interface, both supporting identical end-effectors and diverse tactile sensors (Xense, GelSight Mini, Tac3D, Daimon). \textbf{Middle:} The dataset covers 6 visuo-tactile manipulation patterns, instantiated across 5 semantic scenarios. \textbf{Top-Right:} A scale comparison demonstrates that OmniViTac (21,879 trajectories) significantly exceeds existing visuo-tactile manipulation datasets in tactile-rich data volume. \textbf{Bottom-Right:} The High-quality Data Pipeline ensures reliability through timestamp alignment, visualization, and human-in-the-loop verification.}
\label{fig:dataset_teaser}
\end{figure*}

\section{Related Works}
\subsection{Tactile Sensing and Tactile Representation Learning}
Robotic fine-grained interaction relies on advances in both tactile hardware and representation learning. Existing tactile sensors span multiple physical principles, including capacitive arrays~\cite{xu2024cushsense,xu2016stretch}, piezoresistive sensors~\cite{stassi2014flexible, sundaram2019learning}, and magnetic sensors~\cite{yan2021soft}. Among them, visuo-tactile sensors such as GelSight~\cite{yuan10high} and DIGIT~\cite{lambeta2020digit} have attracted particular attention because they provide high-resolution and robust measurements. By tracking gel deformation with a camera, they reconstruct fine contact geometry and, with marker dots, estimate force and torque.

With such signals available, the focus has shifted from task-specific features to transferable tactile representations. A major line of work learns from paired vision-tactile data through self-supervision or contrastive learning. Sparsh~\cite{higuera2024sparsh} adopts a masked autoencoder, Anytouch~\cite{feng2025anytouch} uses contrastive learning, and UniT~\cite{xu2025unit} introduces a VQGAN-like architecture~\cite{esser2021taming} for compact latent modeling. More recently, tactile sensing has been integrated into MLLMs to align tactile, visual, and textual representations. Representative examples include Binding touch to everything~\cite{yang2024binding}, Octopi~\cite{yu2024octopi}, and VTV-LLM~\cite{xie2025universal}. Octopi-1.5~\cite{yu2025demonstrating} further adds retrieval augmentation to improve predictions for unseen object-tactile pairs. Our study builds on this trend by examining unified visuo-tactile representation learning across four visuo-tactile sensors: GelSight Mini~\cite{yuan10high}, Tac3D~\cite{zhang2022tac3d}, Xense~\cite{xense}, and Daimon~\cite{daimon}.

\subsection{Visuo-Tactile Robotic Manipulation Policies}
Findings from human motor control highlight the importance of tactile feedback for stable manipulation. Prior studies show that cutaneous feedback is essential for maintaining grip stability~\cite{augurelle2003importance} and that the brain uses a unified internal model to predict tactile feedback for both limbs and tools~\cite{kilteni2017sensorimotor, wolpert2001motor}. These results suggest that tactile sensing is fundamental to contact-rich manipulation, where it complements vision under occlusion and provides direct force-related feedback.

Inspired by these insights, robotic manipulation policies increasingly incorporate non-visual inputs. Early visuo-tactile systems such as See-to-Touch~\cite{guzey2024see} and RoboPack~\cite{ai2024robopack} showed that local tactile observations can resolve ambiguities and support fine-grained control when vision is unreliable. Recent work extends this idea to end-to-end manipulation models, including vision-action models such as 3D-ViTac~\cite{huang20243d} and RDP~\cite{xue2025reactive}, as well as vision-language-action models such as VLA-Touch~\cite{bi2025vla}, Tactile-VLA~\cite{huang2025tactile}, and TA-VLA~\cite{zhang2025ta}. However, effectively fusing tactile and visual signals for high-frequency closed-loop control remains difficult. Existing methods often suffer from weak tactile integration or temporal mismatch. Our framework addresses this problem through deeper visuo-tactile fusion, using local tactile feedback to refine global control for robust contact-rich manipulation.

\subsection{Visuo-tactile Manipulation Datasets and Systems}
Early visuo-tactile datasets~\cite{yang2024binding,fu2024touch,yu2024octopi,feng2025anytouch,xie2025universal,yu2025demonstrating} mainly collected static paired observations by pressing handheld tactile sensors against objects, surfaces, or textures. Although useful for multimodal representation learning, they do not provide action trajectories for manipulation. More recent datasets~\cite{li2022see,fang2023rh20t,wu2025freetacman,bi2025vla,cheng2025omnivtla,xu2025exumi,liu2025mla} mount tactile sensors on robotic grippers and record manipulation through UMI~\cite{chi2024universal}, teleoperation, or kinesthetic teaching. These systems capture synchronized vision, touch, and proprioception, enabling visuo-tactile policy learning for contact-rich tasks. Nevertheless, current datasets remain limited in scale, sensor and platform diversity, and temporal alignment between high-frequency tactile streams and visual observations.

Our work addresses these limitations with a multi-platform visuo-tactile manipulation dataset that offers diverse tasks and fully aligned vision, touch, and action sequences, supporting large-scale world-model training for unified visuo-tactile policy learning.
